% Generated draft: requires mathematical and source review.
\documentclass[11pt]{article}
\usepackage[margin=25mm]{geometry}
\usepackage{fontspec}
\setmainfont{Latin Modern Roman}
\usepackage{amsmath,amssymb,mathtools,unicode-math}
\setmathfont{Latin Modern Math}
\usepackage{enumitem,needspace,xurl,hyperref}
\setlength{\parindent}{0pt}
\setlength{\parskip}{5pt}
\setlength{\emergencystretch}{3em}
\newcommand{\N}{\mathbb{N}}
\newcommand{\Z}{\mathbb{Z}}
\newcommand{\Q}{\mathbb{Q}}
\newcommand{\R}{\mathbb{R}}
\newcommand{\C}{\mathbb{C}}
\newcommand{\F}{\mathbb{F}}
\newcommand{\A}{\mathbb{A}}
\newcommand{\PP}{\mathbb{P}}
\newcommand{\E}{\mathbb{E}}
\newcommand{\eps}{\varepsilon}
\newcommand{\dd}{\,\mathrm{d}}
\begin{document}
\section*{UnsolvedMath ID: 20001872}
\textbf{Source problem code:} catalogue entry
\subsection{Definitions and mathematical statement}
Long's source conjecture says every Finsler two-sphere has an elliptic prime closed geodesic. Use a smooth strongly convex Finsler metric, not assumed reversible. Prime means not an iterate of a shorter closed geodesic, and need not mean simple as a curve. On a surface the linearized transverse return map \(P_c\) is in \(\operatorname{Sp}(2,\R)\). The aid uses Long's spectral meaning of elliptic: all eigenvalues lie on the unit circle, including parabolic boundary cases \(|\operatorname{tr}P_c|=2\). Strict ellipticity means \(|\operatorname{tr}P_c|<2\), while hyperbolicity means \(|\operatorname{tr}P_c|>2\). A stricter alternative definition of elliptic would change its residual-branch statement.

The source cites Y. Long--W. Wang [Lon-Wan2008] for at least two elliptic prime closed geodesics on an irreversible Finsler \(S^2\) when the total number of primes is finite. The aid reports the finite-prime branch covered in its intended general setting; the exact extra literature needed beyond the source's irreversible statement is absent. Under that reported finite-branch result, a counterexample is precisely a metric with infinitely many prime closed geodesics, all hyperbolic. It is automatically bumpy: no closed orbit, including any iterate, has transverse multiplier one. This is a conditional reduction using the stated spectral convention, not a proof that such a metric exists.

The source's [Grjuntal1979] example where all closed simple geodesics are hyperbolic does not contradict the conjecture: nonsimple prime orbits may still be elliptic. It also remarks that many primes for some sphere metrics, even some ellipsoids, are nonsimple; no quantitative meaning of most is supplied. The following section heading on complete finite-volume metrics is adjacent extraction context, not a hypothesis of this conjecture.

The candidate compactness proposition assumes \(F_j\to F\) with normalized geodesic vector fields and their linearizations converging on a fixed identified compact unit-tangent model, for example \(C^1\) convergence of those vector fields under radial identifications. Suppose closed \(F_j\)-geodesics \(c_j\), not necessarily prime, satisfy for fixed \(L<\infty\), \(\varepsilon>0\),
\[
 \ell_{F_j}(c_j)\le L,\qquad
 |\operatorname{tr}P_{c_j}|\le2-\varepsilon.
\]
The aid reports that \(F\) then has a strictly elliptic prime closed geodesic. The nonvanishing limiting normalized vector field supplies a positive lower-period bound; without appropriate flow/linearization convergence, an upper length bound alone would not provide the stated conclusion.

Consequently, near any hypothetical Long counterexample, no sequence of strictly elliptic closed geodesics can have both uniformly bounded length and a fixed positive trace gap. Equivalently, for each \(L,\varepsilon\), the two displayed conditions cannot recur along a converging sequence to that counterexample. Length escape and loss of the parabolic gap may occur on different subsequences; the statement does not assert one uniform alternative for every orbit family or exclude all counterexamples. All progress and literature assessments remain attributed unreviewed aid claims.
\subsection{Short English statement}
With spectral ellipticity including parabolic cases, the aid reduces a counterexample to the infinite all-hyperbolic branch. Its compactness result rules out approaching such a metric with elliptic orbits of both bounded length and uniform trace gap. Simple-geodesic hyperbolicity alone does not resolve the conjecture.
\subsection{Sources}
\begin{itemize}
\item Local UnsolvedMath registry, permanent ID \(20001872\), source code \texttt{catalogue entry}; \texttt{aim-geometry-notes.json}, \texttt{notes[209]}, tag \texttt{conjecture}.
\item AIM workshop ``Geodesics,'' recorded source item 2.2, extracted Conjecture 2.2.3 (Long), page 4: \url{https://aimath.org/pastworkshops/geodesicsproblems.pdf}. Source type PDF; statement status \texttt{exact}. Supplied verification identifies page 4. Only the supplied record was used; no PDF was accessed or compiled.
\item Y. Long and W. Wang, source label [Lon-Wan2008], for the finite-prime irreversible case; source label [Grjuntal1979], for a simple-geodesic hyperbolicity example. No authoritative full titles, identifiers, URLs, exact broader finite-prime theorem, or compactness proof artifacts are supplied. Original remarks and literature fields are empty.
\item Research attempt 1: \texttt{reduction}, status at run \texttt{open}, rigor label \texttt{proved}, novelty confidence low; registry status \texttt{partially\_solved}, classification \texttt{PARTIAL-PROGRESS}. Created and updated 14 August 2026. Machine-generated aid, not peer reviewed; independently verify before citation.
\item Ulam AI / UnsolvedMath Contributors; CC BY 4.0.
\end{itemize}
\end{document}
